\section{Geometry and normalization}\label{sec:conventions}

We identify $\SU$ with the unit sphere $S^3$.  Every conjugacy class has a
representative
\[
 g_\theta=\operatorname{diag}(e^{i\theta},e^{-i\theta}),
 \qquad 0\le\theta\le\pi.
\]
Thus $\tfrac12\operatorname{Tr}g_\theta=\cos\theta$.  The Riemannian distance
from the identity is $\theta$.  Normalized Haar
probability has radial part
\begin{equation}\label{eq:haar-radial}
 \dd H(\theta)=\frac{2}{\pi}\sin^2\theta\,\dd\theta.
\end{equation}
Our Laplace--Beltrami operator is nonpositive, and the heat semigroup is
$e^{t\Delta/2}$.  The representation indexed by $n\in\N_0$ has dimension
$n+1$, character
\[
 \chi_n(\theta)=\frac{\sin((n+1)\theta)}{\sin\theta},
\]
and $-\Delta$ eigenvalue $n(n+2)$.  Thus the leading short-time action is
$\theta^2/(2t)$.

\begin{table}[ht]
\centering
\begin{tabular}{lll}
\toprule
Object & convention here & conversion warning \\
\midrule
group metric & unit $S^3$ & Casimir may be scaled by $2$ or $4$ \\
heat equation & $\partial_tu=\tfrac12\Delta u$ & replace $t$ under other generators \\
class angle & $[0,\pi]$ & some physics sources use half the rotation angle \\
Haar measure & probability & volume measure has total mass $2\pi^2$ \\
action & $s_t=-\log p_t$ & additive constants do not affect derivatives \\
\bottomrule
\end{tabular}
\caption{Convention table.  Every validated numerical bound uses the middle
column exactly.}\label{tab:conventions}
\end{table}

These choices are part of every numerical statement.  In the common physics
basis $T^a=\sigma^a/2$, with
$\operatorname{Tr}(T^aT^b)=\tfrac12\delta^{ab}$, one has
$g_\theta=\exp(2i\theta T^3)$.  Consequently a Lie-algebra coordinate written
as $\exp(iX^aT^a)$ has $|X|=2\theta$ in this metric convention.  Constants
must be rescaled before comparison with a convention using $|X|$ itself as
the geodesic distance.
